% ============================================================
% Section 73: 自由电子能态密度（一维/二维/三维）
% ============================================================
\section{固体理论：一维、二维和三维自由电子的能态密度}
\label{sec:free-electron-dos}

\begin{modelbox}[题目：各维度自由电子能态密度的推导]
分别导出一维、二维和三维金属中自由电子的能态密度 $g(E)$。
\end{modelbox}

\begin{tipbox}[考点快速分析]
\begin{itemize}
    \item \textbf{能态密度定义}：$g(E)\,dE$ = 能量在 $[E,E+dE]$ 间隔内的电子状态数（计入自旋简并）
    \item \textbf{核心思路}：利用 $E(\bm{k})$ 的色散关系，将 $k$ 空间的体积/面积/长度元转换为能量间隔
    \item \textbf{通用公式}：$g_d(E)\propto E^{(d-2)/2}$，维度 $d=1,2,3$
    \item \textbf{范霍夫奇点}：一维 DOS 在带底 $E\to0$ 时发散
\end{itemize}
\end{tipbox}

\begin{derivationbox}[标准解题过程]
\textbf{预备知识}

自由电子能量与波矢关系
\begin{equation}
E(k) = \frac{\hbar^2k^2}{2m}
\quad\Longrightarrow\quad
k = \frac{\sqrt{2mE}}{\hbar},\qquad
\dd{k} = \frac{m}{\hbar^2k}\,\dd{E} = \frac{\sqrt{m}}{\hbar\sqrt{2E}}\,\dd{E}.
\end{equation}

计入自旋简并（因子 2），$k$ 空间中每个量子态占据的体积为 $(2\pi/L)^d$。

\textbf{(1) 一维情况}

设一维金属长度为 $L$。$k$ 空间态密度为 $L/(2\pi)$。

能量在 $E\sim E+dE$ 之间的状态，对应 $k$ 空间中 $k\sim k+dk$ 和 $-k\sim -(k+dk)$ 两个区间（正负波矢对称）：
\begin{equation}
g(E)\,dE = 2\times\frac{L}{2\pi}\times 2\,dk = \frac{2L}{\pi}\,dk.
\end{equation}
代入 $\dd{k}=\sqrt{m/(2\hbar^2E)}\,\dd{E}$：
\begin{equation}
g(E) = \frac{2L}{\pi}\cdot\frac{\sqrt{m}}{\hbar\sqrt{2E}}
= \frac{L}{\pi}\sqrt{\frac{2m}{\hbar^2}}\,E^{-1/2}.
\end{equation}

\begin{equation}
\boxed{g_{1D}(E) = \frac{L}{\pi}\Bigl(\frac{2m}{\hbar^2}\Bigr)^{1/2} E^{-1/2}\propto E^{-1/2}.}
\end{equation}

\textbf{关键特征}：当 $E\to0$（带底）时，$g(E)\to\infty$，出现\textbf{范霍夫奇点}。

\textbf{(2) 二维情况}

设二维金属面积为 $S$。$k$ 空间态密度为 $S/(2\pi)^2$。

等能线为半径 $k$ 的圆，能量间隔 $dE$ 对应 $k$ 空间中宽度为 $dk$ 的圆环。圆环面积 $=2\pi k\,dk$：
\begin{equation}
g(E)\,dE = 2\times\frac{S}{(2\pi)^2}\times 2\pi k\,dk = \frac{S}{\pi}\,k\,dk.
\end{equation}

由色散关系 $dE = \frac{\hbar^2}{m}k\,dk$，即 $k\,dk = \frac{m}{\hbar^2}\,dE$。代入得
\begin{equation}
\boxed{g_{2D}(E) = \frac{Sm}{\pi\hbar^2}=\text{常数}.}
\end{equation}

\textbf{关键特征}：二维自由电子的能态密度与能量\textbf{无关}，是一个常数。这是二维电子气最重要的特性之一。

\textbf{(3) 三维情况}

设三维金属体积为 $V$。$k$ 空间态密度为 $V/(2\pi)^3$。

等能面为半径 $k$ 的球面，能量间隔 $dE$ 对应厚度为 $dk$ 的球壳。球壳体积 $=4\pi k^2\,dk$：
\begin{equation}
g(E)\,dE = 2\times\frac{V}{(2\pi)^3}\times 4\pi k^2\,dk = \frac{V}{\pi^2}\,k^2\,dk.
\end{equation}

代入 $k^2=2mE/\hbar^2$ 和 $\dd{k}=\sqrt{m/(2\hbar^2E)}\,\dd{E}$：
\begin{align}
g(E) &= \frac{V}{\pi^2}\cdot\frac{2mE}{\hbar^2}\cdot\frac{\sqrt{m}}{\hbar\sqrt{2E}} \notag\\
&= \frac{V}{2\pi^2}\Bigl(\frac{2m}{\hbar^2}\Bigr)^{3/2} E^{1/2}.
\end{align}

\begin{equation}
\boxed{g_{3D}(E) = \frac{V}{2\pi^2}\Bigl(\frac{2m}{\hbar^2}\Bigr)^{3/2} E^{1/2}\propto E^{1/2}.}
\end{equation}

\textbf{关键特征}：三维 DOS 从带底 $E=0$ 开始按 $\sqrt{E}$ 增长，没有奇点。
\end{derivationbox}

\begin{formulabox}[答案汇总]
\begin{center}
\begin{tabular}{cll}
\toprule
\textbf{维度} & \textbf{能态密度 $g(E)$} & \textbf{能量依赖} \\
\midrule
1D & $\displaystyle g_{1D}(E)=\frac{L}{\pi}\bigl(\frac{2m}{\hbar^2}\bigr)^{1/2}E^{-1/2}$ & $\propto E^{-1/2}$ \\
2D & $\displaystyle g_{2D}(E)=\frac{Sm}{\pi\hbar^2}$ & 常数 \\
3D & $\displaystyle g_{3D}(E)=\frac{V}{2\pi^2}\bigl(\frac{2m}{\hbar^2}\bigr)^{3/2}E^{1/2}$ & $\propto E^{1/2}$ \\
\bottomrule
\end{tabular}
\end{center}
\end{formulabox}

\begin{insightbox}[物理图像]
\begin{itemize}
    \item \textbf{统一规律}：$g_d(E)\propto E^{(d-2)/2}$。维度越高，低能态密度被抑制得越厉害（因为高维 $k$ 空间``更空旷''）。
    \item \textbf{范霍夫奇点}：一维系统在带底和带顶都会出现 DOS 发散，这在碳纳米管、量子线等低维系统中可被实验观测。
    \item \textbf{二维常数 DOS}：石墨烯等二维材料的许多奇异性质（如量子霍尔效应）都与其能量无关的态密度密切相关。
    \item \textbf{三维 $\sqrt{E}$ 律}：这是金属电子热容 $C_V\propto T$、泡利顺磁磁化率与温度无关等性质的根源。
\end{itemize}
\end{insightbox}
